\section{Discussion}
\label{sec:discussion}

\subsection{Decomposition Trade-offs}

The SIG-based decomposition exploits a fundamental property of film scheduling: scenes sharing actors and locations are naturally grouped in practice.
By clustering these scenes into shooting blocks, the upper-level search space reduces from $n!$ scene permutations to $K!$ block permutations ($K \ll n$).

Our results reveal an important trade-off in the number of blocks $K$:
\begin{itemize}
    \item When $K$ is \emph{too small} (e.g., $K=3$ for S15), the permutation space ($3!=6$) is exhaustively searchable, and the block-level ordering provides limited flexibility compared to scene-level search ($15! \approx 1.3 \times 10^{12}$). This explains why GA and PSO outperform BDM on S15 and S20.
    \item When $K$ is \emph{moderate} (e.g., $K=6$ for S25), the decomposition provides a meaningful yet tractable search space ($6!=720$), and the structure-guided ordering yields solutions competitive with scene-level metaheuristics (BDM-AOS within 0.4\% of PSO on S25).
    \item On \textbf{S10}, all methods achieve similar costs (GA: 11{,}780; BDM-AOS: 11{,}965; gap: 1.6\%), validating that the decomposition preserves solution quality on small instances.
\end{itemize}

The key insight is that decomposition becomes increasingly valuable as $n$ grows, because the relative advantage of structured search over random permutation search increases exponentially with problem size.

\subsection{Role of Adaptive Operator Selection}

The Q-learning AOS adapts operator selection to the search state.
On S25, BDM-AOS achieves a cost of 36{,}947 with zero variance across seeds, versus BDM-NoAOS's mean of 37{,}479 (std: 923).
This 1.4\% improvement shows that AOS consistently steers toward the best block ordering, while random operator selection misses it in some runs.
On S10--S20, both variants find identical solutions because the block-permutation spaces ($\leq 24$) are small enough to be fully explored regardless of operator choice.
The most frequently selected operators on S25 are PMX+block\_swap and OX+block\_swap.

\subsection{On the Multi-Objective Formulation}

BDM-AOS uses NSGA-II to jointly optimize cost and makespan.
However, our experiments reveal that the greedy decoder assigns a fixed makespan for each instance (e.g., 7~days for S10, 16~days for S25), regardless of scene ordering.
As a result, the Pareto front collapses to a single point---cost varies with permutation but makespan does not.
This means cost and makespan are \emph{not} conflicting objectives under the current decoder.
To obtain genuine trade-offs, the decoder would need to allow variable daily capacity or idle-day insertion, which we leave for future work.

\subsection{Comparison with Literature}

Our results are consistent with the MSSP literature:
\begin{itemize}
    \item \citet{tekin2023} reported PSO achieving $<1\%$ gaps from optimal on small instances. Our GA and PSO baselines similarly converge to near-identical costs on S10.
    \item \citet{liu2019} reported TSBM gaps of 0.18\% and PSOBM gaps of 7.92\% from MILP optimal. BDM-AOS's gap of 7.4\% from GA on S25 falls in a similar range, though direct comparison is limited by different instance sets.
    \item The zero variance of BDM-AOS across seeds contrasts with GA/PSO variance (e.g., PSO std of 1{,}218 on S25), suggesting the decomposition-guided ordering is deterministic once the best block permutation is found.
\end{itemize}

\subsection{Key Takeaway}

The decomposition approach works best when Louvain produces enough blocks ($K \geq 6$) to give NSGA-II a meaningful search space.
For instances where $K \leq 4$, direct scene-level metaheuristics remain superior.
Future work should target instances with 50+ scenes, where $K$ naturally grows and the $n!/K!$ search-space reduction becomes decisive.
